% Generated draft: requires mathematical and source review.
\documentclass[11pt]{article}
\usepackage[margin=25mm]{geometry}
\usepackage{fontspec}
\setmainfont{Latin Modern Roman}
\usepackage{amsmath,amssymb,mathtools,unicode-math}
\setmathfont{Latin Modern Math}
\usepackage{enumitem,needspace,xurl,hyperref}
\setlength{\parindent}{0pt}
\setlength{\parskip}{5pt}
\setlength{\emergencystretch}{3em}
\newcommand{\N}{\mathbb{N}}
\newcommand{\Z}{\mathbb{Z}}
\newcommand{\Q}{\mathbb{Q}}
\newcommand{\R}{\mathbb{R}}
\newcommand{\C}{\mathbb{C}}
\newcommand{\F}{\mathbb{F}}
\newcommand{\A}{\mathbb{A}}
\newcommand{\PP}{\mathbb{P}}
\newcommand{\E}{\mathbb{E}}
\newcommand{\eps}{\varepsilon}
\newcommand{\dd}{\,\mathrm{d}}
\begin{document}
\section*{UnsolvedMath ID: 2200001}
\textbf{Source problem code:} catalogue entry
\subsection{Definitions and mathematical statement}
Let $N\geq1$, let $x_1,\ldots,x_N\in\R^3$ be distinct fixed charge locations, and let $\xi_1,\ldots,\xi_N$ be their real charge strengths. On $\Omega=\R^3\setminus\{x_1,\ldots,x_N\}$ define
\[
E(x)=\sum_{i=1}^N\frac{\xi_i(x-x_i)}{\|x-x_i\|^3},\qquad
\Phi(x)=\sum_{i=1}^N\frac{\xi_i}{\|x-x_i\|},\qquad E=-\nabla\Phi.
\]
An equilibrium is a point $x\in\Omega$ with $E(x)=0$. The source assumes the set of equilibria is finite and asks whether
\[
\#\{x\in\Omega:E(x)=0\}\leq(N-1)^2.
\]
The count is of distinct points, not multiplicities, and no nondegeneracy assumption is in the conjecture. The source imposes no sign restriction on the strengths; it does not explicitly discuss zero strengths. Charge locations themselves are excluded from the field domain.

The supplied research summary reports a refutation by Arathoon--Ball--Kvalheim, preprint version 1 of July 29, 2026. Its described configuration has $N=5$: unit charges at the vertices of an equilateral triangle and two equal small charges at axial positions $\pm\varepsilon$, with
\[
q_\varepsilon=\frac34\varepsilon^3-\frac5{32}\varepsilon^5.
\]
For sufficiently small positive $\varepsilon$, the summary reports at least $24$ nondegenerate equilibria, exceeding $(5-1)^2=16$, and a generic perturbation of the charge strengths yielding a Morse potential with $k\geq24$ critical points. Nondegeneracy means $\det D^2\Phi\ne0$ at a critical point. The full coordinate normalization and the global finiteness argument are not reproduced in this record and must be checked against the finite-equilibrium hypothesis of the conjecture.

The summary additionally claims its author analytically and numerically rechecked a computational lemma saying that a degree-$4$ limit potential $\Phi_0$ has exactly $21$ nondegenerate critical points. That limit potential and its calculation are not supplied. The $21$-point auxiliary claim is distinct from the $24$-equilibrium claim for the charge configuration. The summary explicitly qualifies the claimed refutation as a very recent, not yet peer-reviewed preprint; its asserted high confidence and independent re-verification belong to the supplied summary, not to this draft.
\subsection{Short English statement}
Maxwell's conjecture in the source bounds a finite set of electrostatic equilibria of $N$ point charges by $(N-1)^2$. The supplied summary reports a July 2026 five-charge construction with at least $24$ equilibria, against the bound $16$, but explicitly identifies the claimed refutation as an unrefereed preprint.
\subsection{Sources}
\begin{itemize}
\item Local UnsolvedMath registry, permanent ID 2200001, exact source code catalogue entry, ``Problems Around Polynomials --- Conjecture 1''; status solved, classification SOLVED-IN-LITERATURE.
\item Shapiro, \emph{Problems Around Polynomials: The Good, The Bad and The Ugly} (2015), source-order item 1, Conjecture 1 as printed, attributed to Maxwell with the source annotations ``seems bad, no tools,'' \url{https://arxiv.org/abs/1503.05295}. The record identifies extraction from public arXiv source TeX, accessed July 29, 2026, and says the source presented this as open.
\item Arathoon, Ball and Kvalheim, arXiv:2607.27197, version 1, July 29, 2026, cited by the supplied research summary for the five-charge refutation. No title or theorem number is supplied. The summary explicitly says this preprint was not yet peer-reviewed, and attributes a high-confidence analytic/numerical recheck and assessment of persistence/transversality arguments to its own author. No such recheck was performed for this draft.
\end{itemize}
\end{document}
